\cotefile{39}
\scanpage{3}
$k$ un \deleted{\unsure{corps commutatif avec}}\deleted{\inserted{\unsure{anneau}}} corps
\[
S_n(k)=k[[t_1,\ldots,t_n]]\qquad(n\ge0;\ S_0(k)=k)
\]
\margin{pour $n$ variable, on a un anneau cosimplicial}
$\mathcal A_n\subset S_n(k)$\qquad sous-anneau \unsure{simplicial|cosimplicial} i.e.

(b) $\forall n\ge0$, $\mathcal A_n$ est un sous-anneau unitaire de $S_n$

(a) On a

$a_1)$ $\mathcal A_n$ est st. par $\unsure{\mathfrak S_n|\mathcal S_n}$ opérant sur $S_n$

$a_2)$ Si $f(t_1,\ldots,t_n)\in\mathcal A_n$, alors $f(t_1,\ldots,t_n)\in\mathcal A_{n+1}$
\inserted{$t_{n+1}$ variable qui ne figure pas}

$a_3)$ Si $f(t_1,\ldots,t_n)\in\mathcal A_n$, alors $f(t_1,t_2,\ldots,\underbrace{t_{n-1},t_{n-1}}_{\text{répété}})\in\mathcal A_{n-1}$
\figure{\begin{tikzpicture}[x=1cm,y=1cm]\draw[->] (.5,2.5)..controls(-.4,2)and(-.3,.5)..(.3,.3)--(1,0);\end{tikzpicture}}
Le plus petit $\mathcal A_*=(\mathcal A_n)_{n\in\mathbb N}$ \deleted{\unsure{pareil}}\inserted{satisfaisant a) b)} est donné par \deleted{\unsure{$\mathcal A_*=P$}}
\[
\mathcal A_n=\operatorname{Im}(\mathbb Z\longrightarrow S_n)
\]
donc l'\deleted{objet}\inserted{anneau} cosimplicial constant de valeur $\operatorname{Im}(\mathbb Z\to k)$,
(qui est $\mathbb Z$ ou un $\mathbb Z/p\mathbb Z$, \deleted{\unsure{si $p>0$}}\inserted{$p$ nb premier} $=\operatorname{car}.k$)

On suppose de plus que

c) Si $f\in\mathcal A_n$\inserted{et} ($f\notin\operatorname{rad}S_n$, i.e. $f$ inv. dans $S_n$, alors $\frac1f\in\unsure{S_n|\mathcal A_n}$.
\[
\Updownarrow
\]
c') $\mathcal A_n$ est un anneau local, et l'inclusion $\mathcal A_n\to S_n$ est un hom local

Le plus petit $\mathcal A$ satisfaisant a) b) c) est formé des
\[
\mathcal A_n=\operatorname{Im}(P\longrightarrow S_n)
\]
où $P$ est le sous-corps premier de $k$ ($\mathbb Q$ en car $0$, $\mathbb F_p$ en car $p>0$).

d) $t_1\in\mathcal A_1\subset\unsure{S|k}[[t_1]]$

Le plus petit $\mathcal A$ satisfaisant à a) b) d) \deleted{\unsure{c)}} est formé des
\[
\mathcal A_n=P[t_1,\ldots,t_n]\deleted{\unsure{_{\mathfrak m_n}}}\qquad\deleted{\unsure{\mathfrak m_n=(}}
\]
le plus petit $\mathcal A$ satisfaisant a) b) c) d) est formé des
\[
\mathcal A_n=P[t_1,\ldots,t_n]_{\mathfrak m_n}\qquad\mathfrak m_n=\sum t_i P[t_1,\ldots,t_n].
\]
\deleted{\unsure{$\Phi$}}
Un hom $S_n\longrightarrow S_m$\inserted{ou $\underbrace{\operatorname{Spf}(S_m)}_{\Phi_m}\unsure{-|\to}\underbrace{\operatorname{Spf}(S_n)}_{\Phi_n}$}, donné par $f_1,\ldots,f_n\in S_m$ est dit un $\mathcal A$-hom si $f_1,\ldots,f_n\in\mathcal A_m$. Il résulte de b) et d) que l'identité $\Phi_n\longrightarrow\Phi_n$, \deleted{\unsure{et les applications}}\inserted{plus gén. les applications $\Phi_m\to\Phi_n$} sont des $\mathcal A$-hom (car $t_i\in\mathcal A_n$ pour $1\le i\le n$). L'analogue de b)
\mspage{2}
